% STATUS: OUTLINE ONLY. Zero points at the midterm. Write it after the
% Evaluation section, because a conclusion written first tends to describe
% the result you hoped for.

\section{Conclusion and Future Work}

% ---------------------------------------------------------------------
% WHAT TO SAY, in about three paragraphs.
%
% 1. What was established. The problem stated as an MDP with a measured
%    action encoding; four baselines; a validated cost assumption on real
%    hardware; and the measured relationship between workload volatility
%    and the room available to a smarter policy.
%
% 2. The main finding. A reactive heuristic is near-optimal under slow
%    drift and falls 41% short of the achievable bound under fast drift,
%    whilst moving nearly three times as much data. The case for learning
%    is therefore conditional on volatility, and that condition is now
%    measured rather than assumed.
%
% 3. What was not achieved. The agent does not yet beat the heuristic.
%    Say it in one sentence, without hedging, and then say what the
%    evidence points at.

% ---------------------------------------------------------------------
\subsection{Future work}
% In the order they should be attempted:
%
% 1. Train at volatility 2 to 4, where there is something to win.
% 2. Address credit assignment: n-step returns, or reward shaping that
%    credits a move with the latency it saves in later windows rather
%    than charging it entirely at the moment it is made.
% 3. Longer training. Run 2 had not converged at 300,000 steps.
% 4. Fine-tune the pre-trained agent on the real cluster. The apparatus
%    exists and the parity constraints are satisfied; the harness is not
%    yet built.
% 5. Choose N by measurement and re-run the cluster experiment.
% 6. Add a cost-model rebalancer in the style of Lion as a stronger
%    baseline than follow traffic.
